\chapter{Unit 3: Ideation and Creative Problem Solving}
\label{chap:unit3_qb}

\section*{Practice Question Set (30 MCQs)}
\addcontentsline{toc}{section}{Practice Question Set (30 MCQs)}

\mcqitem{1}{What is the primary objective of the Ideation stage in Design Thinking?}{Writing production code immediately}{Transitioning from problem definition to generating a broad range of potential solutions}{Finalizing the product marketing budget}{Performing financial audits}

\mcqitem{2}{Which golden rule must be maintained during the initial phase of idea generation?}{Evaluate the financial cost of every idea immediately}{Defer judgment and separate idea generation from idea evaluation}{Allow only senior managers to speak}{Reject any idea that cannot be built within one week}

\mcqitem{3}{In creativity assessment, what does "Fluency" measure?}{How beautifully an idea is drawn}{The total number of ideas generated within a given timeframe}{The grammar and vocabulary used in the pitch}{The technical complexity of the code}

\mcqitem{4}{In creativity assessment, what does "Flexibility" measure?}{The physical flexibility of prototype materials}{The number of different categories or domains represented among the ideas}{The willingness of a team to work overtime}{The adaptability of software servers}

\mcqitem{5}{Who established the foundational rules of modern Brainstorming in 1953?}{Tim Brown}{Alex Osborn}{Karl Duncker}{Steve Jobs}

\mcqitem{6}{Which improv phrase is widely used during brainstorming to build constructively on colleagues' ideas?}{"No, but..."}{"Yes, and..."}{"Why would that work?"}{"That is too expensive"}

\mcqitem{7}{What is "Brainwriting" (such as the 6-3-5 method)?}{Writing software code directly in a notebook}{A silent idea-generation technique where participants write ideas on paper and pass them around to build upon}{Taking notes during a university lecture}{Drawing mind maps individually without sharing}

\mcqitem{8}{In the SCAMPER framework, what does the letter 'S' represent?}{Simplify}{Substitute}{Systematize}{Scale}

\mcqitem{9}{In the SCAMPER framework, replacing paper bus tickets with a mobile QR code scan is an example of:}{Put to another use}{Substitute}{Reverse}{Magnify}

\mcqitem{10}{In SCAMPER, combining a luggage trolley with an electric scooter mechanism is an example of:}{Combine}{Eliminate}{Reverse}{Substitute}

\mcqitem{11}{In SCAMPER, what does 'E' stand for?}{Expand}{Evaluate}{Eliminate}{Energize}

\mcqitem{12}{In SCAMPER, taking a self-checkout system where customers scan goods themselves instead of cashiers scanning them represents:}{Adapt}{Reverse / Rearrange}{Minify}{Substitute}

\mcqitem{13}{How does a Mind Map structure ideas?}{In a linear alphabetical list}{Radially outward from a central problem keyword into branching associative themes}{In a chronological timeline from left to right}{In a 2x2 mathematical matrix}

\mcqitem{14}{What is "Reverse Thinking" (Inversion / Worst Possible Idea)?}{Reading a design brief backwards}{Intentionally brainstorming ways to cause or worsen the problem, then inverting those ideas into solutions}{Copying a competitor's product in reverse}{Writing text from right to left}

\mcqitem{15}{What is the main purpose of "Affinity Diagramming"?}{Grouping large volumes of scattered ideas into natural thematic clusters based on relationships}{Calculating weighted mathematical scores}{Selecting one single winning idea randomly}{Testing prototype durability}

\mcqitem{16}{What is "Concept Bundling"?}{Selling multiple unsold products at a discount}{Packaging several complementary features into a coherent, high-value solution concept}{Archiving rejected ideas into zip files}{Merging two competing design teams}

\mcqitem{17}{Which of the following is a classic two-stage idea selection process?}{Blind voting $\rightarrow$ Random selection}{Initial coarse filter (feasibility/constraints) $\rightarrow$ Deep comparative evaluation matrix}{Manager selection $\rightarrow$ Customer survey}{Lowest cost $\rightarrow$ Immediate launch}

\mcqitem{18}{In a Weighted Screening Matrix, if Criterion A has a weight of 0.6 and Criterion B has a weight of 0.4, an idea scoring 5 on A and 2 on B gets a total score of:}{7.0}{3.8}{4.2}{2.9}

\mcqitem{19}{In an Impact-Effort Matrix, ideas that offer HIGH IMPACT with LOW EFFORT are classified as:}{Major Projects}{Quick Wins}{Thankless Tasks}{Fill-ins}

\mcqitem{20}{In an Impact-Effort Matrix, ideas that require HIGH EFFORT but yield LOW IMPACT should generally be:}{Implemented immediately}{Given the highest budget}{Dropped, avoided, or deprioritized}{Assigned to junior designers}

\mcqitem{21}{What is "Dot Voting" in an ideation workshop?}{A method where participants place colored adhesive dots next to their favorite ideas to visually map group consensus}{Drawing dotted lines across wireframes}{Writing punctuation marks on sticky notes}{Voting by sending text messages}

\mcqitem{22}{In the Double Diamond design model, when does ideation take place?}{The first diamond (Discover)}{The first diamond (Define)}{The opening of the second diamond (Develop / Diverge on solutions)}{After product launch}

\mcqitem{23}{Why is it dangerous to converge on the very first idea that comes to mind during brainstorming?}{The first idea is always illegal}{First ideas are typically obvious, conventional, and lack deeper exploratory value}{First ideas cannot be prototyped}{Brainstorming rules require generating at least 500 ideas}

\mcqitem{24}{What is the role of a facilitator during a team brainstorming session?}{Dictate the best solutions to the team}{Enforce time limits, maintain energy, encourage quiet members, and uphold brainstorming rules}{Criticize unfeasible suggestions immediately}{Write all ideas without allowing team members to participate}

\mcqitem{25}{Which SCAMPER prompt is being used when a smartphone manufacturer removes the physical 3.5mm headphone jack to create more internal battery space?}{Modify}{Eliminate}{Adapt}{Put to another use}

\mcqitem{26}{What does "Originality" represent in idea generation?}{How quickly an idea was conceived}{The statistical rarity and unconventional nature of an idea compared to common standards}{The patent registration number}{The cost required to manufacture the idea}

\mcqitem{27}{When evaluating ideas, why is "Desirability" assessed before "Feasibility"?}{Engineers do not work until marketing finishes}{There is no point engineering a technically feasible system that nobody actually wants or needs}{Feasibility is completely irrelevant in modern design}{Desirability is always cheaper than feasibility}

\mcqitem{28}{What should a design team do with interesting ideas that are not selected for the current MVP?}{Delete them permanently from all records}{Archive them in an "Idea Parking Lot" / backlog for future iterations or reference}{Penalize the team members who created them}{Sell them to competitors}

\mcqitem{29}{What is "Production Blocking" in traditional verbal group brainstorming?}{Running out of sticky notes and markers}{A phenomenon where participants forget or suppress their ideas while waiting for others to stop speaking}{Factory manufacturing machine breakdown}{Managers canceling a brainstorming workshop}

\mcqitem{30}{Which ideation technique is specifically designed to eliminate Production Blocking and evaluation apprehension?}{Open verbal round-robin}{Silent Brainwriting (6-3-5)}{Public debate}{Impromptu pitching}


\newpage
\section*{Answer Key \& Step-by-Step Explanations}
\addcontentsline{toc}{section}{Answer Key \& Explanations}

\begin{answerkeytable}{Unit 3: Quick Answer Key}
\begin{tabular}{c c c c c c c c c c}
\toprule
\textbf{Q1} & \textbf{Q2} & \textbf{Q3} & \textbf{Q4} & \textbf{Q5} & \textbf{Q6} & \textbf{Q7} & \textbf{Q8} & \textbf{Q9} & \textbf{Q10} \\
\textbf{B} & \textbf{B} & \textbf{B} & \textbf{B} & \textbf{B} & \textbf{B} & \textbf{B} & \textbf{B} & \textbf{B} & \textbf{A} \\
\midrule
\textbf{Q11} & \textbf{Q12} & \textbf{Q13} & \textbf{Q14} & \textbf{Q15} & \textbf{Q16} & \textbf{Q17} & \textbf{Q18} & \textbf{Q19} & \textbf{Q20} \\
\textbf{C} & \textbf{B} & \textbf{B} & \textbf{B} & \textbf{A} & \textbf{B} & \textbf{B} & \textbf{B} & \textbf{B} & \textbf{C} \\
\midrule
\textbf{Q21} & \textbf{Q22} & \textbf{Q23} & \textbf{Q24} & \textbf{Q25} & \textbf{Q26} & \textbf{Q27} & \textbf{Q28} & \textbf{Q29} & \textbf{Q30} \\
\textbf{A} & \textbf{C} & \textbf{B} & \textbf{B} & \textbf{B} & \textbf{B} & \textbf{B} & \textbf{B} & \textbf{B} & \textbf{B} \\
\bottomrule
\end{tabular}
\end{answerkeytable}

\begin{expbox}[Detailed Pedagogical Explanations]
\begin{enumerate}[leftmargin=6mm, itemsep=2.5mm]
  \item \textbf{[Q1: Option B]} Ideation transitions from problem framing to creating diverse candidate solutions.
  \item \textbf{[Q2: Option B]} Deferring judgment ensures psychological safety and wide exploratory breadth.
  \item \textbf{[Q3: Option B]} Fluency is the sheer numerical volume of ideas produced.
  \item \textbf{[Q4: Option B]} Flexibility measures the diversity of categories and perspectives represented.
  \item \textbf{[Q5: Option B]} Alex Osborn popularized structured brainstorming in his 1953 book *Applied Imagination*.
  \item \textbf{[Q6: Option B]} "Yes, and..." encourages additive, constructive idea building.
  \item \textbf{[Q7: Option B]} Brainwriting allows silent, parallel idea generation without social intimidation.
  \item \textbf{[Q8: Option B]} In SCAMPER, S stands for Substitute.
  \item \textbf{[Q9: Option B]} Replacing one medium (paper) with another (digital QR code) is Substitution.
  \item \textbf{[Q10: Option A]} Merging two distinct tools into a single entity is Combination.
  \item \textbf{[Q11: Option C]} In SCAMPER, E stands for Eliminate.
  \item \textbf{[Q12: Option B]} Inverting who performs the scanning (customer instead of cashier) is Reverse/Rearrange.
  \item \textbf{[Q13: Option B]} Mind maps radiate organically from a central node through branching concepts.
  \item \textbf{[Q14: Option B]} Reverse thinking explores how to cause a failure and inverts the findings for solutions.
  \item \textbf{[Q15: Option A]} Affinity diagramming synthesizes messy data into thematic clusters.
  \item \textbf{[Q16: Option B]} Concept bundling groups mutually reinforcing features into a single solution package.
  \item \textbf{[Q17: Option B]} A coarse filter weeds out non-starters, followed by detailed multi-criteria matrix scoring.
  \item \textbf{[Q18: Option B]} Calculation: $(0.6 \times 5) + (0.4 \times 2) = 3.0 + 0.8 = 3.8$.
  \item \textbf{[Q19: Option B]} High Impact + Low Effort = Quick Wins (highest initial return for effort).
  \item \textbf{[Q20: Option C]} High Effort + Low Impact ideas are "Money Pits" / "Thankless Tasks" and should be avoided.
  \item \textbf{[Q21: Option A]} Dot voting provides a fast, democratic visual heat-map of team preferences.
  \item \textbf{[Q22: Option C]} In the Double Diamond, Develop is the divergent phase for solution generation.
  \item \textbf{[Q23: Option B]} First ideas are almost always superficial and derived from existing paradigms.
  \item \textbf{[Q24: Option B]} The facilitator maintains process discipline, timekeeping, and psychological safety.
  \item \textbf{[Q25: Option B]} Removing an existing component or feature is an application of "Eliminate".
  \item \textbf{[Q26: Option B]} Originality reflects how novel, unique, and non-standard an idea is.
  \item \textbf{[Q27: Option B]} Validating user need (desirability) ensures technical effort is not wasted.
  \item \textbf{[Q28: Option B]} An Idea Backlog/Parking Lot preserves unselected ideas for future releases.
  \item \textbf{[Q29: Option B]} Production blocking occurs when people must wait their turn to speak and lose ideas.
  \item \textbf{[Q30: Option B]} Brainwriting allows everyone to write simultaneously in silence, eliminating blocking.
\end{enumerate}
\end{expbox}
